\documentclass[11pt,a4paper]{article}

% --- PAQUETES PRINCIPALES PARA OVERLEAF ---
\usepackage[spanish,es-tabla]{babel}
\usepackage[utf8]{inputenc}
\usepackage[T1]{fontenc}
\usepackage{amsmath,amssymb,amsfonts,amsthm,mathtools}
\usepackage{geometry}
\geometry{top=2.5cm, bottom=2.5cm, left=2.5cm, right=2.5cm}
\usepackage{xcolor}
\usepackage{listings}
\usepackage{booktabs}
\usepackage{graphicx}
\usepackage{float}
\usepackage{url}
\usepackage{hyperref}
\usepackage{tcolorbox}
\tcbuselibrary{skins,breakable}
\usepackage{tikz}
\usetikzlibrary{positioning,shapes.geometric,arrows.meta,trees}

% --- PALETA DE COLORES PROFESIONAL / CORPORATIVA ---
\definecolor{corpNavy}{RGB}{18, 52, 86}
\definecolor{corpTeal}{RGB}{0, 128, 128}
\definecolor{corpSlate}{RGB}{74, 85, 104}
\definecolor{lightCard}{RGB}{247, 250, 252}
\definecolor{borderCard}{RGB}{226, 232, 240}
\definecolor{codeBackground}{RGB}{245, 247, 250}
\definecolor{commentGreen}{RGB}{46, 139, 87}
\definecolor{stringCrimson}{RGB}{165, 42, 42}
\definecolor{keywordNavy}{RGB}{25, 25, 112}

\hypersetup{
    colorlinks=true,
    linkcolor=corpNavy,
    citecolor=corpTeal,
    urlcolor=corpTeal,
    pdfencoding=auto
}

% --- CONFIGURACION DE LISTINGS (COMPATIBILIDAD UTF-8 TOTAL) ---
\lstdefinestyle{engineeringStyle}{
    backgroundcolor=\color{codeBackground},
    commentstyle=\color{commentGreen}\itshape,
    keywordstyle=\color{keywordNavy}\bfseries,
    numberstyle=\tiny\color{gray},
    stringstyle=\color{stringCrimson},
    basicstyle=\ttfamily\footnotesize,
    breakatwhitespace=false,
    breaklines=true,
    captionpos=b,
    keepspaces=true,
    numbers=left,
    numbersep=7pt,
    showspaces=false,
    showstringspaces=false,
    showtabs=false,
    tabsize=2,
    frame=single,
    rulecolor=\color{borderCard},
    extendedchars=true,
    literate=
        {á}{{\'a}}1
        {é}{{\'e}}1
        {í}{{\'i}}1
        {ó}{{\'o}}1
        {ú}{{\'u}}1
        {Á}{{\'A}}1
        {É}{{\'E}}1
        {Í}{{\'I}}1
        {Ó}{{\'O}}1
        {Ú}{{\'U}}1
        {ñ}{{\~n}}1
        {Ñ}{{\~N}}1
        {¿}{{?`}}1
        {¡}{{!`}}1
}
\lstset{style=engineeringStyle}

% --- CAJAS DE INFORMACION TECNICA (TCOLORBOX) ---
\newtcolorbox{techcard}[1]{
    colback=lightCard,
    colframe=corpNavy,
    fonttitle=\bfseries\color{white},
    title=#1,
    arc=3mm,
    boxrule=1pt,
    breakable
}

\newtcolorbox{solutioncard}[1]{
    colback=green!5,
    colframe=corpTeal,
    fonttitle=\bfseries\color{white},
    title=#1,
    arc=3mm,
    boxrule=1pt,
    breakable
}

\newtcolorbox{incidentcard}[1]{
    colback=red!5,
    colframe=red!70!black,
    fonttitle=\bfseries\color{white},
    title=#1,
    arc=3mm,
    boxrule=1pt,
    breakable
}

% --- METADATOS DEL INFORME ---
\title{
    \vspace{-1.5cm}
    \textbf{\Large INFORME TÉCNICO DE INGENIERÍA Y DIARIO DE DESARROLLO}\\[0.3cm]
    \textbf{\LARGE Plataforma QuCo: Implementación del Compilador DSL Cuántico, Refactorización de Persistencia y Optimización de Transpilación}\\[0.4cm]
    \large Documento de Seguimiento de Proyecto, Explicación de Arquitectura y Transferencia Técnica
}
\author{\textbf{Equipo de Desarrollo \& Ingeniería Software QuCo}}
\date{\today}

\begin{document}

\maketitle

\begin{abstract}
Este informe técnico documenta de manera detallada las tareas de desarrollo, refactorización y resolución de incidencias realizadas en la plataforma de computación cuántica \textbf{QuCo}. El objetivo de este documento es servir como guía para explicar en reuniones de equipo, revisiones de código y transferencias técnicas qué se ha implementado, cómo está estructurada cada capa tecnológica (Frontend Angular, Backend Spring Boot, Base de Datos MySQL y Entorno de Ejecución Qiskit) y cómo se han resuelto los problemas críticos de negocio detectados (desacoplamiento de compuertas genéricas en base de datos, desbordamiento de memoria por el bug de los 127 cúbits, desvíos de resultados por orden de bits y compilación modular de circuitos complejos).
\end{abstract}

\tableofcontents
\newpage

% ==============================================================================
\section{Resumen Ejecutivo: ¿Qué se ha hecho y por qué?}
% ==============================================================================

En las últimas semanas de trabajo sobre la plataforma \textbf{QuCo}, nos encontramos con varias limitaciones técnicas que impedían avanzar tanto en la experiencia de usuario como en la ejecución real de circuitos cuánticos. Las intervenciones principales han sido:

\begin{enumerate}
    \item \textbf{Implementación de un compilador de lenguaje DSL cuántico en el cliente}:
    Permite describir circuitos complejos como fórmulas algebraicas (por ejemplo, el algoritmo de Grover $g\_4\_2\_7 = h(4) \cdot [O(4, 2) \cdot O(4, 7) \cdot D(4)]^{\{2\}} \cdot M(4, '+') $), analizando la sintaxis mediante un analizador léxico (\textit{lexer}) y sintáctico (\textit{parser}) que genera automáticamente la rejilla visual y el código Python Qiskit.
    
    \item \textbf{Desacoplamiento total de las compuertas genéricas mediante Base de Datos}:
    Anteriormente, el frontend tenía nombres de compuertas cableados en código (\textit{hardcoded}). Ahora, la base de datos MySQL y el backend Spring Boot son la única fuente de verdad, gestionando un nuevo atributo booleano \texttt{isGeneric} que categoriza las compuertas automáticamente.
    
    \item \textbf{Corrección del problema de sobreasignación de 127 cúbits}:
    Al conectar la plataforma con procesadores reales de IBM Quantum (como \texttt{ibm\_brisbane}), el sistema asignaba 127 cúbits a todos los circuitos. Esto colapsaba los simuladores de vectores de estado. Se implementó una lógica de dimensionamiento dinámico según los cúbits efectivamente en uso.
    
    \item \textbf{Corrección de la discrepancia de orden de bits (\textit{Endianness}) en el oráculo de Grover}:
    Se corrigió el generador de oráculos para respetar la ordenación little-endian de Qiskit frente al formato big-endian de interfaces tipo Quirk, garantizando que el algoritmo encuentre con alta probabilidad los estados diana y no estados erráticos.
\end{enumerate}

\begin{figure}[H]
\centering
\resizebox{0.95\textwidth}{!}{%
\begin{tikzpicture}[
    node distance=1.2cm and 1.4cm,
    layer/.style={rectangle, draw=corpNavy, fill=lightCard, rounded corners, text width=3.6cm, text centered, minimum height=1.2cm, font=\small},
    arrow/.style={-{Stealth[scale=1.2]}, corpTeal, very thick}
]
    \node[layer] (front) {\textbf{Frontend (Angular)}\\ Editor Visual, DSL Parser, Transpilador en Cliente};
    \node[layer, right=of front] (back) {\textbf{Backend (Spring Boot)}\\ REST API, CircuitEditorService, Inicializador Idempotente};
    \node[layer, right=of back] (db) {\textbf{Base de Datos (MySQL)}\\ Tabla \texttt{qiskit\_code}\\ Campo \texttt{is\_generic}};
    \node[layer, below=of back] (exec) {\textbf{Runtime Cuántico}\\ Qiskit Python, AerSimulator, IBM Quantum Platform};

    \draw[arrow] (front) -- node[above, font=\footnotesize] {JSON / DTOs} (back);
    \draw[arrow] (back) -- node[above, font=\footnotesize] {JPA / JDBC} (db);
    \draw[arrow] (back) -- node[right, font=\footnotesize] {Ejecución Código} (exec);
    \draw[arrow] (exec) -| node[above, font=\footnotesize] {Histogramas} (front);
\end{tikzpicture}%
}
\caption{Arquitectura funcional de extremo a extremo de la solución implementada.}
\label{fig:arch_overview}
\end{figure}

% ==============================================================================
\section{Implementación de la Persistencia y Backend (Spring Boot + MySQL)}
% ==============================================================================

\subsection{El Problema de Partida: Compuertas Cableadas}
En versiones anteriores de QuCo:
\begin{itemize}
    \item El cliente web clasificaba las compuertas cuánticas mediante comprobaciones de cadenas literales (\texttt{if (gate.name === 'InitialH' || gate.name === 'Diffuser')}).
    \item No era posible añadir una nueva compuerta genérica sin modificar el código fuente de TypeScript y volver a desplegar la aplicación.
    \item Existía desincronización entre el código Qiskit que generaba el frontend y el que guardaba el backend.
\end{itemize}

\subsection{Paso 1: Migración en Base de Datos MySQL}
Se alteró la tabla \texttt{qiskit\_code} en el esquema \texttt{quco} para incorporar un identificador de tipo booleano:
\begin{lstlisting}[language=SQL, caption={Script DDL ejecutado en MySQL}]
ALTER TABLE qiskit_code ADD COLUMN is_generic BOOLEAN DEFAULT FALSE;
\end{lstlisting}

\subsection{Paso 2: Modelado JPA en Java (QiskitCode.java)}
En el backend Spring Boot (\path{tp3/src/main/java/.../QiskitCode.java}), se mapeó la columna para permitir que Hibernate la procese y Jackson la serialice hacia el cliente:

\begin{lstlisting}[language=Java, caption={Definicion de atributo en la entidad QiskitCode.java}]
@Column(name = "is_generic")
@JsonProperty("isGeneric")
private Boolean isGeneric = false;

public Boolean getIsGeneric() {
    return isGeneric;
}

public void setIsGeneric(Boolean isGeneric) {
    this.isGeneric = isGeneric;
}
\end{lstlisting}

\subsection{Paso 3: Inicialización Automática Idempotente (GenericGatesInitializer)}
Para evitar que un despliegue en un entorno nuevo o una base de datos limpia requiera introducir manualmente las compuertas genéricas, se desarrolló un componente \texttt{CommandLineRunner} que se ejecuta en el arranque del servidor Spring Boot:

\begin{lstlisting}[language=Java, caption={Servicio de Inicializacion Idempotente GenericGatesInitializer.java}]
@Component
public class GenericGatesInitializer implements CommandLineRunner {

    @Autowired
    private QiskitCodeRepository qiskitCodeRepo;

    @Override
    public void run(String... args) {
        seedGateIfMissing("InitialH", getInitialHPythonCode(), true);
        seedGateIfMissing("Oracle", getOraclePythonCode(), true);
        seedGateIfMissing("Diffuser", getDiffuserPythonCode(), true);
        seedGateIfMissing("MCP", getMcpPythonCode(), true);
        seedGateIfMissing("Measure", getMeasurePythonCode(), true);
        seedGateIfMissing("0", getSubcircuito0PythonCode(), true);
    }

    private void seedGateIfMissing(String name, String pythonCode, boolean isGeneric) {
        Optional<QiskitCode> existing = qiskitCodeRepo.findByName(name);
        if (existing.isEmpty()) {
            QiskitCode code = new QiskitCode();
            code.setName(name);
            code.setCode(pythonCode);
            code.setIsGeneric(isGeneric);
            qiskitCodeRepo.save(code);
        } else if (!Boolean.TRUE.equals(existing.get().getIsGeneric())) {
            QiskitCode code = existing.get();
            code.setIsGeneric(isGeneric);
            qiskitCodeRepo.save(code);
        }
    }
}
\end{lstlisting}

\begin{solutioncard}{Beneficio Técnico Conseguido}
El arranque del backend es totalmente autosuficiente: si las compuertas no existen, las crea; si existen pero no tenían marcado el flag \texttt{is\_generic}, las actualiza sin borrar datos del usuario.
\end{solutioncard}

% ==============================================================================
\section{Implementación del Compilador DSL en el Frontend Angular}
% ==============================================================================

Para soportar la especificación formal del documento conceptual donde los circuitos se definen mediante expresiones como:
\begin{equation}
    g\_4\_2\_7 = h(4) \cdot [O(4, 2) \cdot O(4, 7) \cdot D(4)]^{\{2\}} \cdot M(4, \text{'+'})
\end{equation}
se construyó un motor de compilación modular en TypeScript dentro del directorio \path{front/src/app/circuit-editor/dsl/}.

\subsection{Módulo 1: Analizador Léxico (\texttt{QuantumDslLexer})}
El lexer se encarga de descomponer el flujo de texto en una secuencia de fichas (\textit{tokens}) tipadas, reconociendo operadores especiales de composición ($\cdot$, $*$, $\bullet$), repetición ($\{k\}$), bloques ($[ ]$), paréntesis y cadenas literales.

\begin{table}[H]
\centering
\caption{Tokens principales procesados por el QuantumDslLexer.}
\label{tab:tokens_dsl}
\resizebox{0.92\textwidth}{!}{%
\begin{tabular}{@{}lll@{}}
\toprule
\textbf{Token} & \textbf{Símbolo / Patrón} & \textbf{Uso en el Lenguaje} \\ \midrule
\texttt{IDENTIFIER} & \texttt{[a-zA-Z\_][a-zA-Z0-9\_]*} & Nombre de compuerta o variable (\texttt{h}, \texttt{O}, \texttt{Diffuser}) \\
\texttt{NUMBER} & \texttt{[0-9]+} & Cantidad de cúbits o repeticiones (\texttt{4}, \texttt{2}, \texttt{7}) \\
\texttt{DOT} & \texttt{·}, \texttt{*}, \texttt{.} & Operador de concatenación temporal secuencial \\
\texttt{LBRACE / RBRACE} & \texttt{\{ \}} & Factor de repetición de bucle (\texttt{\{2\}}) \\
\texttt{LBRACKET / RBRACKET} & \texttt{[ ]} & Agrupación de subcircuito o lista de argumentos \\
\texttt{EQUALS} & \texttt{=} & Asignación de nombre al circuito \\ \bottomrule
\end{tabular}%
}
\end{table}

\subsection{Módulo 2: Analizador Sintáctico (\texttt{QuantumDslParser})}
El parser utiliza una estrategia de descenso recursivo sin retroceso (\textit{backtracking}) para transformar la lista de tokens en un Árbol de Sintaxis Abstracta (\textit{Abstract Syntax Tree}, AST).

\begin{figure}[H]
\centering
\resizebox{\textwidth}{!}{%
\begin{tikzpicture}[
    level 1/.style={sibling distance=5.2cm, level distance=1.4cm},
    level 2/.style={sibling distance=2.6cm, level distance=1.4cm},
    level 3/.style={sibling distance=2.1cm, level distance=1.4cm},
    nodeStyle/.style={rectangle, draw=corpNavy, fill=lightCard, rounded corners, font=\small, align=center, inner sep=5pt}
]
    \node[nodeStyle] {\textbf{AssignmentNode}\\ name: \texttt{"g\_4\_2\_7"}}
        child { node[nodeStyle] {\textbf{SequenceNode}\\ Operador $\cdot$}
            child { node[nodeStyle] {\textbf{CallNode}\\ \texttt{h(4)}} }
            child { node[nodeStyle] {\textbf{RepeatNode}\\ count: \texttt{2}}
                child { node[nodeStyle] {\textbf{SequenceNode}\\ Subcircuito}
                    child { node[nodeStyle] {\textbf{CallNode}\\ \texttt{O(4, 2)}} }
                    child { node[nodeStyle] {\textbf{CallNode}\\ \texttt{O(4, 7)}} }
                    child { node[nodeStyle] {\textbf{CallNode}\\ \texttt{D(4)}} }
                }
            }
            child { node[nodeStyle] {\textbf{CallNode}\\ \texttt{M(4, '+')}} }
        };
\end{tikzpicture}%
}
\caption{Representación en memoria del AST generado para la expresión de Grover.}
\label{fig:ast_diagram}
\end{figure}

\subsection{Módulo 3: Motor de Evaluación y Generación de Código (\texttt{QuantumDslEngine})}
El motor de evaluación toma el AST y realiza dos tareas simultáneas:
\begin{enumerate}
    \item \textbf{Colocación en el Editor Visual de QuCo}:
    Convierte cada compuerta en una instancia \texttt{GatePlacement} asignándole índice de cúbit y columna de tiempo en el grid de QuCo. Soporta dos modalidades:
    \begin{itemize}
        \item \textbf{Modo Expandido}: Descompone los oráculos y difusores en sus compuertas atómicas ($H$, $X$, $\text{MCP}$) paso a paso.
        \item \textbf{Modo Modular}: Empaqueta oráculos y difusores como cajas compuestas en una sola columna.
    \end{itemize}
    \item \textbf{Generación de Código Qiskit Dual}:
    Genera tanto el script Python plano directo como la versión modular basada en llamadas compuestas (\texttt{circuit.compose(...)}).
\end{enumerate}

\subsection{Módulo 4: Descompilador Inverso \texorpdfstring{(Python $\to$ DSL)}{(Python -> DSL)}}
Se implementó la función \texttt{decompileGenericGatePythonCode(pythonCode)} para permitir que cuando el usuario inspecciona o edita una compuerta en el editor de código, QuCo sea capaz de inferir la fórmula equivalente en el DSL. El descompilador reconoce patrones como:
\begin{itemize}
    \item Bucles \texttt{for \_ in range(k):} $\longrightarrow$ Expresión repetida \texttt{[...]}\{k\}.
    \item Bloques \texttt{for i in range(n): U.x(i)} con condiciones $\longrightarrow$ \texttt{Binary(n, m)}.
    \item Bloques \texttt{target / controls / U.mcp} $\longrightarrow$ \texttt{mcp(n, pi, n-1)}.
\end{itemize}

% ==============================================================================
\section{Resolución de Incidencias Críticas en Producción}
% ==============================================================================

\subsection{Incidencia 1: El Bug de los 127 Cúbits}

\begin{incidentcard}{Descripción de la Incidencia}
Al seleccionar procesadores cuánticos físicos de IBM Quantum (como \texttt{ibm\_brisbane}), la plataforma forzaba un valor estático:
\begin{lstlisting}[language=Python]
qubits = 127
outputQubits = 127
\end{lstlisting}
Al enviar este circuito al simulador local de vector de estado (\texttt{Statevector}), el proceso de Python intentaba reservar un vector complejo de dimensión $2^{127}$, provocando un fallo de memoria instantáneo (\textit{Out Of Memory}).
\end{incidentcard}

\subsubsection{Causa Raíz Identificada}
En el servicio backend (\path{CircuitEditorService.java}) y en el componente cliente (\path{circuit-editor.component.ts}), el conteo de cúbits usaba el tamaño del chip (\texttt{circuit.getQubits().size()}) en vez de los cúbits activos con compuertas.

\subsubsection{Solución de Ingeniería Aplicada}
Se implementó un algoritmo de inspección que escanea todas las compuertas activas del circuito y calcula el índice máximo ocupado:
\begin{equation}
    N_{\text{activos}} = \max \big( \{ \text{gate.qubit} \mid \text{gate} \in \text{circuit.gates} \} \big) + 1
\end{equation}

\begin{lstlisting}[language=Java, caption={Correccion en CircuitEditorService.java}]
// 1. Deteccion del indice maximo ocupado
int maxQubitIndex = -1;
for (EdGate gate : circuit.getGates()) {
    if (gate.getQubit() > maxQubitIndex) {
        maxQubitIndex = gate.getQubit();
    }
}

// 2. Si hay compuertas en q[0]..q[3], maxQubitIndex = 3 -> usedQubitsCount = 4
int usedQubitsCount = (maxQubitIndex >= 0) ? (maxQubitIndex + 1) : 4;

finalCode = finalCode.replace("#QUBITS#", String.valueOf(usedQubitsCount));
finalCode = finalCode.replace("#OUTPUT_QUBITS#", String.valueOf(usedQubitsCount));
\end{lstlisting}

En la interfaz de usuario de QuCo, se aplicó la misma regla:
\begin{lstlisting}[language=HTML, caption={Correccion en circuit-editor.component.ts}]
const maxQubitIndex = this.circuit.gates.reduce((max, g) => Math.max(max, g.qubit), -1);
const totalQubits = maxQubitIndex >= 0 ? (maxQubitIndex + 1) : Math.min(this.circuit?.qubits.length || 4, 4);
\end{lstlisting}

\begin{solutioncard}{Resultado}
Circuitos de 4 cúbits diseñados en el entorno de \texttt{ibm\_brisbane} ahora generan \texttt{qubits = 4}, permitiendo simularlos en milisegundos en máquinas locales y ejecutarlos mediante mapeo de topología física en IBM Quantum.
\end{solutioncard}

\subsection{Incidencia 2: Inversión de Resultados por Orden de Bits (Endianness)}

\begin{incidentcard}{Descripción de la Incidencia}
Al configurar un oráculo para buscar los números $2$ y $8$, la simulación devolvía distribuciones anómalas con picos en $1$ y $4$.
\end{incidentcard}

\subsubsection{Causa Raíz}
En el entorno gráfico Quirk, los cúbits se ordenan de mayor a menor significancia (big-endian), mientras que Qiskit utiliza little-endian (el bit menos significativo va en \texttt{q[0]}). Una conversión de cadena invertida producía la transposición de bits:
\begin{equation}
    |0010\rangle_2 \; (2) \xrightarrow{\text{lectura invertida}} |0100\rangle_2 \; (4)
\end{equation}

\subsubsection{Solución de Ingeniería Aplicada}
Se eliminó la manipulación manual de cadenas y se implementó la extracción directa de bits mediante el operador de desplazamiento (\textit{shift}) aritmético:
\begin{lstlisting}[language=Python, caption={Extraccion directa de bits sin riesgo de inversion}]
# Bit 'i' corresponde exactamente al qubit q[i] de Qiskit
for i in range(n):
    if ((m >> i) & 1) == 0:
        U.x(i)
\end{lstlisting}
Tras esta modificación, la búsqueda de $2$ y $7$ en 4 cúbits produce de forma consistente una probabilidad de acierto del **$96.4\%$** en el simulador.

% ==============================================================================
\section{Guía de Operación y Puesta en Marcha}
% ==============================================================================

Esta sección resume los comandos y comprobaciones necesarias para levantar el entorno de desarrollo y probar las nuevas funcionalidades.

\subsection{Arranque de los Servicios}

\subsubsection{1. Base de Datos MySQL}
Asegurarse de que el servidor MySQL en el puerto 3306 esté activo y que el usuario de la aplicación tenga permisos sobre el esquema \texttt{quco}:
\begin{lstlisting}[language=SQL]
SELECT table_name, column_name, data_type 
FROM information_schema.columns 
WHERE table_name = 'qiskit_code' AND column_name = 'is_generic';
\end{lstlisting}

\subsubsection{2. Backend Spring Boot (Puerto 8080)}
Navegar al directorio del backend e iniciar el servicio:
\begin{lstlisting}[language=bash]
cd tp3
mvn spring-boot:run
\end{lstlisting}
El log del servidor mostrará la inicialización automática de compuertas:
\begin{verbatim}
[INFO] GenericGatesInitializer: Seeding default generic gates
       (is_generic = true)...
[INFO] InitialH, Oracle, Diffuser, MCP, Measure verified.
\end{verbatim}

\subsubsection{3. Frontend Angular (Puerto 4200)}
En otra terminal, iniciar el servidor de desarrollo web:
\begin{lstlisting}[language=bash]
cd front
npm start
\end{lstlisting}
Acceder en el navegador a \url{http://localhost:4200/circuit-editor}.

\subsection{Resolución de Conflictos de Puertos (Puerto 8080)}
Si el puerto 8080 se encuentra ocupado por un proceso previo colgado, se puede liberar de forma limpia mediante:
\begin{lstlisting}[language=bash]
# Consultar el PID del proceso en el puerto 8080
lsof -ti :8080

# Detener el proceso
lsof -ti :8080 | xargs kill -9
\end{lstlisting}

% ==============================================================================
\section{Check-list de Demostración Técnica para Reuniones}
% ==============================================================================

Para presentar y justificar este trabajo ante el equipo técnico, se recomienda seguir este flujo de demostración:

\begin{enumerate}
    \item \textbf{Demostración de Persistencia Desacoplada}:
    Abrir el inspector de red en el navegador y mostrar la llamada GET a \texttt{/qiskit/getCodes}. Observar que cada compuerta incluye el campo \texttt{"isGeneric": true/false}. Modificar una compuerta en la base de datos y comprobar que el frontend reacciona dinámicamente sin recompilar código.
    
    \item \textbf{Demostración del Compilador DSL}:
    En el editor de circuitos de QuCo, abrir la consola del DSL e introducir:
    \begin{verbatim}
    g_4_2_7 = h(4) · [O(4, 2) · O(4, 7) · D(4)]{2} · M(4, '+')
    \end{verbatim}
    Pulsar \textbf{Compilar}: observar cómo se pueblan instantáneamente las 17 columnas de la rejilla con sus compuertas $H$, $X$ y $\text{MCP}$.
    
    \item \textbf{Demostración de Ejecución Cuántica}:
    Ejecutar el circuito en el simulador Aer y mostrar el histograma resultante: los estados decimales $2$ ($|0010\rangle$) y $7$ ($|0111\rangle$) aparecen con una concentración de probabilidad superior al $90\%$.
    
    \item \textbf{Demostración del Dimensionamiento de Cúbits}:
    Seleccionar el procesador \texttt{ibm\_brisbane} (127 cúbits). Pulsar en ver código Qiskit y verificar que el encabezado genera:
    \begin{lstlisting}[language=Python]
    qubits = 4
    outputQubits = 4
    \end{lstlisting}
    demostrando que los 123 cúbits restantes no sobrecargan la memoria del simulador.
\end{enumerate}

% ==============================================================================
\section{Conclusiones y Próximos Pasos}
% ==============================================================================

El conjunto de cambios implementados ha transformado a QuCo de una herramienta de edición visual manual a una plataforma con soporte para lenguajes específicos de dominio y ejecución en hardware real:
\begin{itemize}
    \item Se ha resuelto la deuda técnica histórica de compuertas cableadas en el cliente.
    \item Se ha establecido un marco formal para componer algoritmos cuánticos de forma algebraica.
    \item Se ha garantizado la estabilidad de la plataforma al eliminar los cuellos de botella de memoria en simulación.
\end{itemize}

\textbf{Próximas mejoras planificadas}:
\begin{enumerate}
    \item Optimización de circuitos en el DSL: agregar un pase de simplificación algebraica que cancele compuertas redundantes consecutivas ($H \cdot H = I$, $X \cdot X = I$).
    \item Soporte para oráculos basados en funciones aritméticas booleanas y grafos de conectividad física.
\end{enumerate}

\end{document}
